% Capitulo: Fundamentos que não envelhecem
\chapter{Fundamentos que não envelhecem}
\label{cap:fundamentos}

\rascunho{
  \textbf{Ideia do capítulo.} O que vale estudar porque não muda com a troca de ferramenta.
  \begin{itemize}
    \item Modelagem de dados.
    \item Estatística básica e pensamento crítico com números.
    \item Qualidade de dados.
    \item Comunicação: escrever, apresentar, contar o que o dado mostra.
    \item Conhecer o negócio: aqui só a apresentação, o assunto ganha o capítulo~\ref{cap:negocio} inteiro.
  \end{itemize}
}

\section{Modelagem}
\label{sec:fund-modelagem}

\pendente{escrever}

\section{Estatística e pensamento crítico}
\label{sec:fund-estatistica}

\pendente{escrever}

\section{Qualidade de dados}
\label{sec:fund-qualidade}

\pendente{escrever}

\section{Comunicação}
\label{sec:fund-comunicacao}

\pendente{escrever}

\section{Conhecer o negócio}
\label{sec:fund-negocio}

\pendente{escrever uma ponte curta para o capítulo~\ref{cap:negocio}}
